% !TeX root = ../papers/nc-p3.tex
\section{真の二重根を与えるKummer四次曲面}
\label{sec:kummer-ramification}
$z\lambda(z^2-1)(\lambda^2z-1)(\lambda^2-1)(\lambda^2z^2-1)\ne0$とする。
この節の射影空間は$A_0$の係数$[a:b:c:d]$の空間である。
$D_0=ad-bc$と書き、節\ref{sec:klein-elliptic-completion}の判別式から
\[
 F=\lambda^2(z^2-1)^2D_0^2\,\delta
\]
を作る。これは全射影空間上の斉次四次式である。
\begin{theorem}\label{thm:kummer-ramification}
上の倍率条件を満たす全ての幾何点$(z,\lambda)$で、
$F=0$は幾何学的に整な四次曲面であり、特異点は相異なる十六の通常二重点で尽きる。
この曲面は六点で分岐する種数2曲線のJacobianの符号商と双有理である。
$P,Q$がともに真の二次式で重根を持つ許容窓は
$x=-m/(2p)$, $y=-k/(2q)$から一意に復元される。
この曲面を$(z,\lambda)$とともに動かした既約な族には、全幾何点が標準
Koszul左右AS正則で反対階数16となる非空開集合がある。
これは$(z,\lambda)=(3,2)$の曲面上の開集合も含む。開集合を縮めると
最小周期は4であり、四つの二十点配置それぞれの射影安定群は指定した四元群である。
\end{theorem}
\begin{proof}
$F=Aa^2+Ba+C$と展開し、$[1:0:0:0]$から$[b:c:d]$の平面に射影する。
$A,B,C$は順に二、三、四次である。次の六直線を取る。
\[
 b-c,\quad b-\lambda d,\quad -c+\lambda z d,\quad z^2 b-c,\quad
 \lambda c-d,\quad \lambda z b-d .
\]
$B^2-4AC$はこれらの積に$16(\lambda^2-1)(z-1)^2(\lambda^2z^2-1)$を
掛けたものである。$e=z\lambda+1/\lambda$とすると六直線の係数点は
双対平面の対称行列
\[
 N=\begin{pmatrix}
 2&1+z^2&e\\
 1+z^2&2z^2&ze\\
 e&ze&2z
 \end{pmatrix}
\]
の二次曲線にある。行列式は
$2z(\lambda^2-1)(z-1)^2(\lambda^2z^2-1)/\lambda^2$で非零、
かつ$A$のHesse行列は$-2\lambda^2\operatorname{adj}(N)$である。
従ってこれらは$A=0$への接線である。
各係数点の零座標と残り二座標の比を比べると六点は相異なり、滑らかな
双対二次曲線上の三点は共線でない。従って六直線の二本ずつの交点は十五個で、
いずれも$A=0$の外にある。以上の係数恒等式は
\path{computations/verify_kummer_discriminant.py}で検算する。

$A\ne0$では平方完成により単純分岐直線上で滑らか、二線交点上に各一個の
通常二重点を持つ。六直線の積は平方でないから関数体は平面の二次拡大である。
$A=0$と$B=0$の共通点は六接点に限られ、そこでは
$d(B^2-4AC)=-4C\,dA\ne0$なので$C\ne0$となる。
よって有限な$a$を持つ追加の特異点も、共通因子$\gcd(A,B,C)$もない。
射影中心$[1:0:0:0]$の接錐は非退化な$A$であり、第十六の通常二重点となる。
超曲面の整性は原始な二次式と関数体の既約性から従う。

代数閉体上で滑らかな二次曲線を$\mathbb P^1$と同定し、
六つの接点を分岐点とした二重被覆$C_2\to\mathbb P^1$を取る。
射影平面を$\operatorname{Sym}^2\mathbb P^1$として、二次曲線を二重点の軌跡とすれば、
指定点を含む次数2因子の直線がその接線である。
$C_2:v^2=f(t)$の二点について$v_1v_2$の平方は、両点が指定分岐点を含む
六つの直線式の積に定数と有理関数の平方を掛けたものである。
従って平面の上のこの二次拡大は、$\operatorname{Sym}^2 C_2$を同時の
超楕円反転で割った関数体に等しい。一般の次数2の線束はRiemann--Rochにより
有効因子を一意に持つ。標準因子を引くAbel写像で$\operatorname{Sym}^2 C_2$を
Jacobianと双有理に同定すると、同時反転は符号反転となる。

具体的に$(z,\lambda)=(3,2)$とし、二行列を
\[
 A_0=\begin{pmatrix}36/17&0\\9&4\end{pmatrix},\qquad
 A_1=\begin{pmatrix}36/19&36/37\\36/5&-12\end{pmatrix}
\]
とする。$p,q,x,y,D_0,\det A_1$は全て非零で$\delta=0$である。
両隣接切断は階数6、反対切断は階数16である。
四次中心は文字次数順に$(3,2,0,0)$次元で、既定の四中心切断で割った商は
素数1009において全始点・経路長13まで$(1+t+t^2+t^3)^4$の係数を与える。
既に正則な一般点を持つ五次元の二次補完族上にあり、次数ごとの標準Hilbert関数の
下界とこの上界から、定理\ref{thm:klein-rational-period-four}の
自由表示および左右Koszul列の議論により全次数で正則となる。

開集合についても理由を述べる。$pqD_0\ne0$における曲面族は、
上の幾何学的に整な四次の一般ファイバーを持ち、四次式に基底だけの因子がなく
整である。標準次元を持つ経路長4,5の近傍で、周囲の十五次元成分の一般点の
四次中心五次元は特殊化先でも下界となる。上の標本も五次元なので、
近傍で可換方程式は定階数となり、四切断を延ばせる。
次数13までの商次元の上界と反対階数16を保つ開条件を課せば、
同じ上下界と自由表示の証明がこの一つの開集合の全幾何点に適用できる。
また素数34381でこの有理標本の四位置の二十点は完全分裂し、四回の核の移動は恒等、
配置間の射影写像の個数行列は$4I_4$である。各配置の三点は非共線で五点は非共面であり、
系\ref{cor:klein-four-configurations}の射影枠の特殊化を適用できる。
四元群はもとから作用するので、開近傍で四つの安定群が確定し、位置の非同値性から
最小周期4が従う。
\end{proof}
